\section[Gradient-Square Closure]{Gradient-Square Closure and Cone Geometry}
\label{sec:gradient-square-closure}

After the two-generator case, a natural candidate for the next closure
is
\[
 R=|x|^2,\qquad S=|\nabla F|^2,\qquad
 \mathbb R[R,F,S].
\]
The generator degrees are \((2,3,4)\) for cubics and \((2,4,6)\)
for quartics. In odd degree, closure forces degree three and identifies
the possible cubic types. In even degree, we record the closure
identities, a regular-link equation and a rectangular matrix example.
Here closure always means that this prescribed
polynomial algebra is Laplacian.

\subsection{Odd-degree rigidity and cubic closure}
\label{sec:gradient-square-rigidity}

The identity
\[
 \cL(F)=S\Delta F-\tfrac12\Gamma(F,S)
\]
relates the closure to minimality. In odd degree, the critical set
determines which of the following two rigidity arguments applies.

\begin{theorem}[Odd-degree rigidity]\label{thm:odd}
Let \(N\geq3\), and let \(F\not\equiv0\) be a homogeneous polynomial of
odd degree \(e\geq3\) on \(\R^N\). Suppose \(\R[R,F,S]\) is a
Laplacian algebra and \(\nabla F(p)=0\) at some \(p\ne0\).
Then \(e=3\).
\end{theorem}

\Cref{sec:odd-rigidity-proof} proves the theorem by degree counting
and Hessian comparisons at a nonzero critical point and a spherical
maximum, forcing $(e-2)(e-3)=0$.

With trivial critical set, the gradient square is radial:

\begin{theorem}[Odd closure with trivial critical set]
\label{thm:odd-isolated-closure}
Let \(N\geq3\), and let \(F\not\equiv0\) be homogeneous of odd degree
\(e\geq3\). If \(\R[R,F,S]\) is Laplacian and
\(\{\nabla F=0\}=\{0\}\), then
\[
 S=cR^{e-1}\qquad(c>0).
\]
Consequently \(e=3\), and, up to an orthogonal change of variables
and a nonzero constant factor, \(F\) is a Cartan cubic in one of the
dimensions \(N=5,8,14,26\).
\end{theorem}

A spherical gradient-flow first integral and incompatible Hessian
traces force $S$ to be radial. Odd closure also gives harmonicity;
the polynomial eikonal classification \cite[Theorem~1]{TkachevIsoparametric14}
and Cartan's cubic classification \cite{Cartan39} then apply.
See \cref{sec:odd-isolated-proof}.
Thus both cases force degree three; the higher-degree pullbacks in
\cref{sec:quartic-application} instead use \(Q=|P|^2\) as a generator.

The restriction \(N\geq3\) is essential. On \(\R^2\),
\(F(x,y)=\operatorname{Re}(x+iy)^e\) is harmonic and satisfies
\(S=e^2R^{e-1}\) for every positive integer \(e\).

\begin{lemma}[Three-generator cubic closure]
\label{lem:cubic-gradient-closure}
Let \(F\) be a nonzero homogeneous cubic on \(\R^N\).  Suppose
\(\R[R,F,S]\) is a Laplacian algebra, where \(S=|\nabla F|^2\).
There are constants \(\alpha,\beta,\gamma,\delta,\varepsilon\) such that
\begin{equation}\label{eq:cubic-five-constants}
\begin{aligned}
 \Delta F&=0,& \Delta S&=\alpha R,\\
 \Gamma(F,S)&=\beta RF,&
 \Gamma(S,S)&=\gamma R^3+\delta F^2+\varepsilon RS.
\end{aligned}
\end{equation}
Moreover, \(\alpha>0\) and
\(\cL(F)=-\beta RF/2\), so the regular zero set is minimal.
\end{lemma}

\begin{proof}
The generators have degrees \(2,3,4\).  Their homogeneous monomials
of degrees \(1,2,5,6\) are, respectively, none, \(R\), \(RF\), and
\(R^3,F^2,RS\).  Laplacian closure and the product rule
\eqref{eq:Gamma-polarization}
give \eqref{eq:cubic-five-constants}.  Bochner's identity
\cite[Proposition~1.47, equation~(1.236), p.~61]{ColdingMinicozzi11} gives
\(\alpha R=2|\Hess F|^2\); the Hessian of a nonzero cubic is not
identically zero, so \(\alpha>0\).  The expression for \(\cL(F)\)
follows from harmonicity.
\end{proof}

At a positive maximum \(a=F(x)\) on the unit sphere, Euler's
identity gives \(\nabla F(x)=3ax\) and
\(\Hess F(x)x=6ax\), hence \(\beta=108a^2>0\).
This gives a harmonic radial eigencubic. Its normalization and
classification follow; variational properties are treated in
\cref{sec:gradient-square-rank-two}.

\subsubsection{Normalized closure and the critical set}

\begin{proposition}[Normalized closure and the critical set]
\label{prop:cubic-singular-normalization}
Assume the hypotheses of \cref{lem:cubic-gradient-closure} and that
\(\nabla F(p)=0\) for some \(p\neq0\).  After multiplying \(F\)
by a positive constant, the closure becomes
\begin{equation}\label{eq:cubic-normalized-closure}
\begin{aligned}
 \Delta F&=0,& \Delta S&=\alpha R,&
 \alpha&=\frac{2(N+3)}3,\\
 \Gamma(F,S)&=4RF,&
 \Gamma(S,S)&=4RS+12F^2.
\end{aligned}
\end{equation}
There is an integer \(k\geq1\) such that
\begin{equation}\label{eq:cubic-dimension-quantization}
 N=6k-3,\qquad \alpha=4k.
\end{equation}
At every nonzero critical point, \(\Hess F\) has \(k\) positive
and \(k\) negative eigenvalues, all of absolute value \(|p|\), and
\(N-2k\) zero eigenvalues.
The punctured critical set \(Z^*=\{x\ne0:\nabla F(x)=0\}\) is a
real-analytic submanifold of dimension \(N-2k=4k-3\), with
\(T_pZ^*=\ker\Hess F(p)\).  In orthonormal coordinates
\(z=(u,v,w)\), where \(u,v\in\R^k\) are the positive and negative
Hessian directions and \(w\in\R^{N-2k}\),
\[
 F(p+z)=\frac{|p|}{2}\bigl(|u|^2-|v|^2\bigr)+O(|z|^3).
\]
\end{proposition}

\begin{proof}
At the given nonzero critical point, Bochner's identity makes
the trace-free Hessian nonzero, while the differentiated Euler
identity places \(p\) in its kernel. Hence \(N\geq3\).
The normalization \eqref{eq:normalizedodd}, established in the
proof of \cref{thm:odd} in \cref{sec:odd-rigidity-proof}, applies
with \(e=3\). Together with the moment identity
\eqref{eq:alphadim}, it gives \eqref{eq:cubic-normalized-closure}.

Let \(p\ne0\) be any critical point and put \(H=\Hess F(p)\).
The cubic Taylor expansion is
\[
 F(p+y)=\tfrac12\langle Hy,y\rangle+F(y),\qquad
 S(p+y)=\langle H^2y,y\rangle+O(|y|^3).
\]
Comparing quadratic terms in \(\Gamma(F,S)=4RF\) gives
\(2\langle H^3y,y\rangle=2|p|^2\langle Hy,y\rangle\),
and hence \(H^3=|p|^2H\). Bochner's identity implies \(H\ne0\),
and harmonicity gives \(\tr H=0\). Thus the nonzero eigenvalues
are \(\pm|p|\), with equal multiplicities, say \(k\).
Bochner's identity and \eqref{eq:cubic-normalized-closure} imply \(2k|p|^2=\tr(H^2)=\tfrac12\alpha|p|^2\).
Thus \(\alpha=4k\) and \(N=6k-3\), proving
\eqref{eq:cubic-dimension-quantization}.

The implicit function argument in \cref{sec:cubic-local-geometry}
proves the asserted description of \(Z^*\). The Hessian eigenvalues
give the displayed Taylor expansion; \cref{sec:tangent-instability-transfer}
then identifies the quadratic tangent and the regular graphical convergence.
\end{proof}

\subsubsection{The cubic classification within this closure}
\begin{proposition}[Cubic types under gradient-square closure]
\label{prop:cubic-gradient-classification}
\label{rem:cubic-radial-eigencubic}
Under the hypotheses of \cref{prop:cubic-singular-normalization},
the possible parameters are
\begin{equation}\label{eq:cubic-realized-dimensions}
 k\in\{1,2,3,5\},\qquad N\in\{3,9,15,27\}.
\end{equation}
These four cases are realized, respectively, by \(xyz\) on \(\R^3\)
and
\[
 F(A)=\tfrac13\operatorname{Re}\tr(A^3)
 \quad\text{on }\operatorname{Herm}_4(\mathbb F)_0,
 \qquad \mathbb F=\R,\mathbb C,\mathbb H.
\]
In particular, for an irreducible defining polynomial only the
dimensions \(9,15,27\) occur.
\end{proposition}

\begin{proof}
We apply the existing radial-eigencubic classification, with its
normalization made explicit. By
\cref{prop:cubic-singular-normalization},
\[
 \cL(F)=-2RF,\qquad
 \tr\bigl((\Hess F)^2\bigr)=2kR.
\]
The Hessian-trace identity in
\cite[Lemma~5.1(i)]{TkachevClassification10} gives the cubic trace
coefficient \(-(N+6)\lambda-6\beta_2\). Substituting
\(\lambda=-2\), \(\beta_2=2k\), and \(N=6k-3\) gives \(\tr\bigl((\Hess F)^3\bigr)=6F\).
The type parameters \((n_1,n_2)\) are the characteristic dimensions
of \cite[Definition~6.7.1, Proposition~6.11.2, and
Corollary~6.11.1]{NadirashviliTkachevVladut14}.
The dimension and cubic Hessian-trace coefficient give
\[
 n_1=1+\frac{6}{3(-2)}=0,
 \qquad n_2=\frac{N+1}{2}=3k-1.
\]
Under the additional quadratic Hessian-trace identity, a radial
eigencubic is either exceptional or one of the four Hurwitz eigencubics
of Clifford type \cite[Corollary~6.11.2 and Section~6.11.5]{NadirashviliTkachevVladut14}.
For the exceptional branch, the \(n_1=0\) classification in
\cite[Proposition~6.1]{TkachevClassification10} gives exactly the
three possible types \((n_1,n_2)=(0,5),\ (0,8),\ (0,14)\).
Among the four Clifford-type eigencubics, \(n_1=0\) occurs only for
the real product cubic in dimension three; the four algebras are listed
as \emph{mutants} in \cite[Section~6.11.5]{NadirashviliTkachevVladut14}.
This proves \eqref{eq:cubic-realized-dimensions}. The three
exceptional cases are realized by the rank-four trace models listed in
\cite[Section~6.11.6(iii)]{NadirashviliTkachevVladut14};
their closure is also calculated directly in
\cref{cor:quaternionic-trace-cubic}. The dimension-three model is
reducible, so it is excluded by irreducibility.
\end{proof}

The general classification of radial eigencubics remains incomplete
\cite[Section~2.2]{FoxTkachev24}. Within this subclass,
\cref{thm:cubic-rank-two-existence} selects only the dimension-27
quaternionic trace cone. General cubic algebras begin in
\cref{sec:cubic-algebraic-theory}.


\subsection{Even-degree closure and a regular-link identity}
\label{sec:even-gradient-structure}

Even-degree gradient-square closure need not imply harmonicity.
We record its identities and a regular-link consequence, leaving
further geometry and variational properties aside.

\begin{proposition}[The even-degree closure and its signs]
\label{prop:even-gradient-closure}
Let \(F\) be a real irreducible homogeneous polynomial of degree
\(d=2m\ge4\) on \(\R^N\), with a nonempty minimal regular zero set.
Suppose \(\R[R,F,S]\) is Laplacian, where
\(R=|x|^2\) and \(S=|\nabla F|^2\).
There are constants \(c,\alpha,\beta,\lambda,\mu\) such that
\begin{equation}\label{eq:even-rfs-closure}
\begin{aligned}
 \Delta F&=cR^{m-1},&
 \Delta S&=\alpha R^{d-2}+\beta R^{m-2}F,\\
 \cL(F)&=(\lambda R^{d-2}+\mu R^{m-2}F)F,\\
 \Gamma(F,S)&=2cR^{m-1}S
       -2\lambda R^{d-2}F-2\mu R^{m-2}F^2.
\end{aligned}
\end{equation}
Moreover, \(\alpha>0\) and \(\lambda<0\).
\end{proposition}

\begin{proof}
The degrees of \(R,F,S\) are \(2,d,2d-2\).
The homogeneous pieces of degrees \(d-2\) and \(2d-4\) are
therefore spanned by \(R^{m-1}\) and by
\(R^{d-2},R^{m-2}F\), respectively.
Minimality and \cref{prop:divisibility} make \(\cL(F)/F\)
a polynomial of degree \(2d-4\); division inside a Laplacian
algebra, \cref{prop:MR-finite-and-division}, places it in the same
homogeneous piece. The identity
\(\cL(F)=S\Delta F-\Gamma(F,S)/2\) proves the mixed entry.

Bochner's identity and \(\Gamma(R,F)=2dF\) give
\begin{equation}\label{eq:even-bochner}
 |\Hess F|^2=\frac\alpha2R^{d-2}
       +\left(\frac\beta2-cd(d-2)\right)R^{m-2}F.
\end{equation}
At a nonzero regular zero this shows \(\alpha\ge0\).
The polynomial \(F\) changes sign across that zero.
If \(\alpha=0\), nonnegativity of the left side would force
\(\beta/2-cd(d-2)=0\), and hence \(\Hess F=0\), contrary
to \(d\ge4\). Thus \(\alpha>0\).

Let \(a>0\) and \(b_0<0\) be the maximum and minimum of
\(F\) on the unit sphere. At either extremum, with value \(t\),
Euler's identity gives
\[
 \nabla F=dt x,\qquad S=d^2t^2,\qquad
 \nabla S=2(d-1)Sx.
\]
Substitution in the mixed entry, followed by division by \(2t\),
gives
\[
 d^3(d-1)t^2+(\mu-cd^2)t+\lambda=0.
\]
The two distinct values \(a,b_0\) are its roots. Consequently
\(\lambda=d^3(d-1)ab_0<0\).
\end{proof}

\begin{remark}[Harmonic normalization]
\label{rem:even-harmonic-suspension}
With \(a=c/[d(N+d-2)]\), the polynomial \(H=F-aR^m\)
is harmonic, and
\[
 |\nabla H|^2=S-2ad^2R^{m-1}F+a^2d^2R^{d-1}.
\]
Thus \(\R[R,H,|\nabla H|^2]=\R[R,F,S]\).
This change preserves the algebra, but generally changes the zero set
and the critical set; it does not reduce the geometric classification
to the harmonic case.

\end{remark}

\begin{proposition}[An isoparametric function on the regular link]
\label{prop:even-singular-link-structure}
Assume \cref{prop:even-gradient-closure} and that $F$ has a nonzero
critical point. Put $\ell=d-1$ and
$h=(\alpha-2c^2)/(-\lambda)$. On
$M=\{F=0,\ S>0\}\cap\Sph^{N-1}$, the function
$u=-2\ell^2S|_M/\lambda$ satisfies $0<u\le1$ and
\begin{equation}\label{eq:even-normalized-link}
 |\nabla_Mu|^2=4\ell^2u(1-u),\qquad
 \Delta_Mu=2\ell^2(h-1)-2\ell(N+2\ell-3)u.
\end{equation}
Thus $u$ is isoparametric in the intrinsic sense that its gradient
square and Laplacian depend only on its value.
\end{proposition}

\begin{proof}
Let $p\ne0$ be critical, $R_0=R(p)$ and $H=\Hess F(p)$.
Euler's identity gives $F(p)=0$. Comparing quadratic terms at $p$
in \eqref{eq:even-rfs-closure} and using \eqref{eq:even-bochner} gives
\begin{equation}\label{eq:even-critical-hessian}
 H(2H^2-2cR_0^{m-1}H+\lambda R_0^{d-2}\Id)=0,
 \quad \tr H=cR_0^{m-1},\quad
 \tr(H^2)=\tfrac\alpha2R_0^{d-2}.
\end{equation}
The normalized nonzero eigenvalues are among
$\eta_\pm=(c\pm\sqrt{c^2-2\lambda})/2$, with opposite signs.
Both occur: $H\ne0$ since $\alpha>0$, and if only $\eta_+$
occurred with multiplicity $a\ge1$, its trace would give
$(a-1)\eta_+=\eta_-$, a contradiction. The other case is identical.

Closure and degree counting give
\begin{align*}
 \Gamma(S,S)={}&D_0R^{2d-3}+D_1R^{3m-3}F
 +D_2R^{d-3}F^2+D_3R^{m-3}F^3\\
 &+D_4R^{d-2}S+D_5R^{m-2}FS,
\end{align*}
with the $D_3$ term omitted for $d=4$.
At $p$ its constant term gives $D_0=0$, and its quadratic terms give
$4H^4=D_1R_0^{3m-3}H/2+D_4R_0^{d-2}H^2$.
For each of the two roots, this reads $4\eta^3=D_1/2+D_4\eta$.
Using $\eta^2=c\eta-\lambda/2$ yields
\begin{equation}\label{eq:even-critical-gamma}
 D_0=0,\qquad D_1=-4c\lambda,\qquad D_4=4c^2-2\lambda.
\end{equation}

On $C_{\reg}=\{F=0,S>0\}$ take $\nu=\nabla F/\sqrt S$.
The mixed entry in \eqref{eq:even-rfs-closure} and
\eqref{eq:even-critical-gamma} imply
\[
 \nu(S)=2cR^{m-1}\sqrt S,\qquad
 |\Hess F\,\nu|^2=(c^2-\lambda/2)R^{d-2}.
\]
Subtracting $\nu(S)^2$ from $\Gamma(S,S)$ gives the tangential
gradient square. Differentiating $\Gamma(F,S)$ in direction $\nu$,
and using $\nu(R)=0$ on $C$, gives
$\Hess S(\nu,\nu)=(2c^2-\lambda)R^{d-2}$.
Since $C_{\reg}$ is minimal,
$\Delta_CS=\Delta S-\Hess S(\nu,\nu)$, and hence
\begin{equation}\label{eq:even-cone-structure}
 |\nabla_CS|^2=-2\lambda R^{d-2}S,\qquad
 \Delta_CS=-(h-1)\lambda R^{d-2}.
\end{equation}
Finally $S$ has degree $2\ell$. Writing $s=S|_M$ and separating
radial derivatives in the cone metric gives
\[
 |\nabla_Ms|^2=-2\lambda s-4\ell^2s^2,\qquad
 \Delta_Ms=-(h-1)\lambda-2\ell(N+2\ell-3)s.
\]
Normalization proves \eqref{eq:even-normalized-link}; its first
identity and $s>0$ imply $0<u\le1$.
\end{proof}

\subsection{Quartic closure and a rectangular matrix family}
\label{sec:quartic-gradient-closure}

The following example realizes generator degrees $(2,4,6)$ by
rectangular matrices, with minimality selecting a single row dimension.

\begin{proposition}[The rectangular matrix family]
\label{ex:rectangular-gradient-closure}
\label{prop:rectangular-minimal-selection}
Let \(p\ge3\), give \(M_{p,3}(\R)\) its Frobenius inner product,
and put
\begin{equation}\label{eq:rectangular-quartic-definition}
 B=X^{\mathsf T}X,\qquad P_j=\tr(B^j),\qquad
 R=P_1,\qquad F=2P_2-R^2,\qquad S=|\nabla F|^2.
\end{equation}
Then \(\cA_R(F)=\R[R,F,S]\), with three algebraically independent
generators of degrees \((2,4,6)\). Its complete closure table is
\begin{equation}\label{eq:rectangular-quartic-laplacians}
 \Delta R=6p,\qquad \Delta F=4(6-p)R,\qquad
 \Delta S=32(p+6)R^2+448F,
\end{equation}
\begin{equation}\label{eq:rectangular-quartic-gamma}
 \begin{aligned}
 \Gamma(R,R)&=4R,\quad \Gamma(R,F)=8F,\quad \Gamma(R,S)=12S,\\
 \Gamma(F,F)&=S,\\
 \Gamma(F,S)&=8RS+128(R^2+F)F,\\
 \Gamma(S,S)&=192(R^2+F)S+1024R^3F+2048RF^2.
 \end{aligned}
\end{equation}
The regular zero set is minimal if and only if \(p=5\). In that
case \(N=15\), \(F\) is nonharmonic, and
\begin{equation}\label{eq:rectangular-quartic-minimality}
 \Delta F=4R,\qquad \cL(F)=-64(R^2+F)F.
\end{equation}
\end{proposition}

\begin{proof}
Differentiation of
\(B=X^{\mathsf T}X\) gives
\begin{equation}\label{eq:rectangular-trace-differentiation}
 \nabla P_j=2jXB^{j-1},\qquad
 \Gamma(P_i,P_j)=4ijP_{i+j-1}.
\end{equation}
For the required Laplacians, the product rule gives
\[
 \Delta P_1=6p,\qquad
 \Delta P_2=4(p+4)P_1,\qquad
 \Delta P_3=6\bigl((p+5)P_2+P_1^2\bigr).
\]
These formulas, derived in \cref{app:rectangular-trace-details},
together with the product rule give
\begin{equation}\label{eq:rectangular-S-traces}
 \nabla F=4X(2B-R\Id),\qquad
 S=64P_3-64RP_2+16R^3
  =64P_3-32RF-16R^3,
\end{equation}
and \eqref{eq:rectangular-quartic-laplacians}.

The characteristic equation of the $3\times3$ matrix $B$ reduces
$P_4,P_5$ to $R,F,S$. The resulting formulas in
\cref{app:rectangular-trace-details} give the last two gradient
entries of \eqref{eq:rectangular-quartic-gamma}; Euler's identity
supplies its first row. Thus \cref{prop:MR-generator-test} proves closure;
conversely every Laplacian algebra containing \(R,F\) contains
\(S=\Gamma(F,F)\), proving the equality with \(\cA_R(F)\).
On diagonal rectangular matrices with three distinct positive squared
entries \(t_1,t_2,t_3\), the Jacobian of
\((P_1,P_2,P_3)\) is a nonzero Vandermonde determinant times six.
The triangular change in \eqref{eq:rectangular-S-traces} is invertible,
so \(R,F,S\) are algebraically independent.

Finally the displayed table gives
\[
 \cL(F)=4(5-p)RS-64(R^2+F)F.
\]
At a matrix with \(B=\operatorname{diag}(4,1,1)\), one has
\(F=0\), \(R=6\), and \(S=768>0\). Thus the regular zero
set can be minimal only for \(p=5\); for this value the curvature
identity \eqref{eq:div-normalized-gradient} proves minimality.
The Laplacian in \eqref{eq:rectangular-quartic-laplacians} proves
the last assertion. For comparison, the member with \(p=6\) is
harmonic and has the same gradient pairings, but its regular zero
set is not minimal.
\end{proof}



\subsection{The boundary of gradient-square closure}

The quadratic compositions of the next section have a different
three-generator structure:

\begin{proposition}[An independent quartic generator]
\label{prop:quartic-composition-not-rfs}
Let \(R,T,F\) be algebraically independent homogeneous polynomials
of degrees \(2,4,4\), with \(R=|x|^2\), and put
\(S=|\nabla F|^2\). Suppose, for constants \(a,b,h\) with \(b>0\),
\[
 S=4R(aF+bT),\qquad
 \Delta S=8(N+8)(aF+bT)+8(a^2+bh)R^2.
\]
Then \(\R[R,F,S]\) is not Laplacian, although
\[
 \R[R,F,S][R^{-1}]=\R[R,T,F][R^{-1}].
\]
\end{proposition}

\begin{proof}
The degree-four part of \(\R[R,F,S]\), whose displayed generators
have degrees \(2,4,6\), is contained in
\(\operatorname{span}\{R^2,F\}\). Algebraic independence places
\(T\) outside this span. Its coefficient in \(\Delta S\) is
\(8(N+8)b\ne0\), so \(\Delta S\notin\R[R,F,S]\).
After adjoining \(R^{-1}\), however,
\(T=b^{-1}(S/(4R)-aF)\), which proves both inclusions in the
displayed equality.
\end{proof}

The orthogonal quartic blocks in \cref{sec:two-quadratic} satisfy
these identities by \eqref{bh:reduced-table}, with
\(\cA_R(F)=\R[R,T,F]\). Thus the \((2,4,6)\) and \((2,4,4)\)
closures differ as polynomial algebras despite agreeing after localization.
